%----------------------------------------------------------------------------
\chapter{Kutatási terv és értékelési módszertan}
%----------------------------------------------------------------------------

A dolgozat gyakorlati részében az kerül vizsgálatra, hogy milyen az elkészített AI alapú szoftver mennyire járul hozzá a tanulók témaspecifikus szókincsfejlődéséhez, beszédkészségéhez, hibafelismeréséhez és tanulási motivációjához a hagyományos tanítási módszerhez viszonyítva. Mindezt egy rövidtávú kontrollcsoportos pilot vizsgálat során derítjük ki.\cite{vanTeijlingen2001pilot}

A vizsgálat egy 40 perces angol nyelvű tanóra keretében valósul meg. Ennek témája a technológiahasználat és az online biztonság, amelyben kitérünk a kártékony szoftverek típusaira is. A tananyag sok új, fontos fogalmat dolgoz fel, mint például: virus, worm, trojan horse, ransomware, personal information, password, suspicious link. A választott téma alkalmas a szókincsfejlesztés, a szövegértés és beszédalapú gyakorlás vizsgálatára is.

A kutatásban összesen 16 tanuló vesz részt, akik egy magántanár diákjai. A tanulók életkora 12 és 16 év között mozog, nyelvi szintjük pedig A2 és B2 közötti. A nemi eloszlás azonos, tehát 8 fiú és 8 lány van, minden csoportban 4-4 közülük. A résztvevők két 8 fős csoportra osztódnak, amiben vegyesen van gyengébb és erősebb szintű tanuló is. A kontrollcsoport hagyományos tanítási módszerrel dolgozza fel a témát, míg a kísérleti csoport AI alapú alkalmazás segítségével tanul.

Mivel a kutatásban kiskorú tanulók vesznek részt, az adatok kezelése anonim módon történik. A neveik nem jelennek meg eredményekben, sem adatbázisban, helyette K1-K8 kódokkal jelöljük őket, míg az AI alapú alkalmazást használó csoportot A1-A8 kóddal. 

A vizsgálatnak a mérésére használunk előtesztet, utótesztet \cite{campbell1963experimental}, beszédkészség mérést és kérdőívet. Az előteszt célja a tanulók kiindulási tudásszintjének mérése, míg az utóteszt a tanulási folyamat után elért eredmények mérésére szolgál. A beszédkészség vizsgálata rövid, irányított beszélgetések alapján történik, míg a kérdőív a tanulói élményt, motivációt és az alkalmazás használhatóságát méri.

\section{A kutatás célja}

A kutatás célja annak vizsgálata, hogy az AI alapú nyelvoktatás milyen szinten támogatja a beszédfejlesztést, szókincs gyarapítást, motivációt, a tanulók rövid távú tanulási eredményei megvizsgálásával, úgy hogy egy konkrét angol nyelvi témát dolgoznak fel. 

A kutatás, mivel két különböző módszert vizsgál, összehasonlító jellegű. Az egyik csoport hagyományosabb, tanári irányítással tanul, amelyben prezentáció, vocabulary feladat és beszélgetés az alapja a lecke elsajátításának. A másik csoport ugyanazt a témát boncolgatja, viszont az alkalmazás keretein belül. Itt a leckét elolvashatja, kijelölheti az ismeretlen kifejezéseket, beszélget az AI aszisztenssel a témáról, amely vizsgálja közben a hibáit, végül pedig a hibákat átismétli.

A kutatás rövidtávú vizsgálatként értelmezhető. Ennek megfelelően az eredmények célja nem általános érvényű következtetéseket von le, hanem annak előzetes vizsgálatára szolgál, hogy az alkalmazás alkalmas lehet oktatási környezetben, és milyen lehetőségeket nyújt a hagyományos módszerekhez képest.

\section{Kutatási kérdések és hipotézisek}

A kutatás során négy darab kérdés tevődik fel, amelyhez kapcsolódik négy hipotézisem.

A kérdések a következők:
\begin{quote}
\begin{enumerate}
    \item Milyen mértékben fejlődik a tanulók témaspecifikus angol szókincse az AI alapú alkalmazás használata után?
    \item Van-e különbség a hagyományos tanítási módszerrel és az AI alapú alkalmazással tanuló csoport eredményei között?
    \item Milyen hatással van az AI alapú alkalmazás használata a tanulók beszédkészségére egy adott témáról folytatott beszélgetés során?
    \item Hogyan értékelik a tanulók az AI alapú alkalmazás használhatóságát, érthetőségét, motiváló hatását?
\end{enumerate}
\end{quote}

A kutatási kérdésekre a következő hipotézisek fogalmazhatók meg:
\begin{quote}
\begin{enumerate}
    \item Az AI-alapú alkalmazást használó tanulók témaspecifikus szókincse fejlődik az előteszt és az utóteszt eredményei alapján.
    \item Az AI-alapú alkalmazást használó csoport fejlődése nagyobb mértékű lesz, mint a hagyományos módszerrel tanuló kontrollcsoporté.
    
    \item Az AI-alapú alkalmazást használó tanulók beszédkészsége javul az adott témáról folytatott irányított beszélgetés során.
    
    \item A tanulók pozitívan értékelik az AI-alapú alkalmazás használhatóságát, érthetőségét és motiváló hatását.
\end{enumerate}
\end{quote}

\section{A tanóra témája és menete}
A tanórák 40 percből állnak, amelyeken kétféleképpen fogjuk megtanítani a diákokat a technológiahasználatról, amelyben sok új kifejezést is tanulhatnak és a beszédkészségüket is növelhetik. Ebben a témában fogunk tesztelni.

A kontrollcsoport esetében az óra hagyományosan, tanári magyarázattal és prezentációval történik. Miután ez megvalósult, vocabulary feladatokat kezdenek oldani, utána pedig brainstorming résszel folytatják, amelyben a tanulók a témához kapcsolódó kérdésekre válaszolnak.

A kísérleti csoport ugyanazt a témát dolgozza fel, az applikáció segítségével. Az applikációban megtalálja ugyanazt a tananyagot szöveg formájában, amelyet meg is hallgathat. Itt kijelöli az ismeretlen kifejezéseket, amelyeket később tanulhat. Ezek után AI által vezetett beszélgetés történik, ahol a témáról beszélgetnek. Itt hibavizsgálat folyik a háttérben beszélgetés közben, és ezután a megállapított hibákat újra begyakorolhatja, és elsajátíthatja helyesen a kifejezéseket.

A két tanulási folyamat eltérő, de ugyanarra a tananyagra készíti fel a tanulókat. Ez lehetővé teszi, hogy egy adott téma tanulási módszerét vizsgáljuk.

\section{Mérőeszközök}
A kutatás során több mérőeszközt is használok, annak érdekében, hogy a tanulási eredményt több szempontból is megvizsgáljuk. Az értékelés kiterjed szókincsfejlődésre, beszédkészségre, tanulói motivációra is. Ennek megfelelően különböző méréseket veszek igénybe.

Elsősorban írásbeli tesztekkel mérek. Ez előteszt és utóteszt formájában valósul meg, \cite{campbell1963experimental} ahol a tanulók kiindulási tudásszintjének felmérése történik meg a tanóra előtt és a tanóra utáni megszerzett tudást is felméri. A feladatok között szerepel a szó-jelentés párosítás, mondatkiegészítés, rövid válaszok írása kérdésekre. 

A tanulói motiváció és élmény mérése kérdőív kitöltésével történik. A kérdőívben 1 és 5 között kell pontozni az állításokat, amelyek vizsgálják, hogy a tanulók mennyire találják hasznosnak, érdekesnek, motiválónak a tanulási módszert.

Az alkalmazással tanuló csoportoknál adatforrásként szolgál az alkalmazásban rögzített tanulási adatok, mint az ismeretlen szavak száma, vocabulary, helytelen mondatok. Ezek az adatok azt is lehetővé teszik, hogy vizsgáljuk, milyen mértékben használták az alkalmazás funkcióit.

\section{Adatelemzési módszerek}
Az adatok elemzése kvantitatív és kvalitatív szempontokkal történik. A kvantitatív elemzés méri, hogy a tanulók teljesítménye mennyire változik az előteszt és utóteszt között, és azt is hogy a kontrollcsoport és a kísérleti csoport eredményei között megfigyelhető lehet különbség vagy sem.

Elsőszámú mérés a leíró statisztikai mutatók kiszámítása. Itt kiszámoljuk az átlagot, szórást, minimum, maximum értékeket előteszt és utóteszt közötti fejlődésben. Ugyanez alkalmazható a beszédkészség felmérésében is, ahol az előtesztre és utótesztre is kap a tanártól pontot.

A csoportokon belüli fejlődés vizsgálatára páros t-próba alkalmazható, \cite{field2013statistics} ahol a tanulók előtesztjei és utótesztjeinek eredményét hasonlítjuk. Ezzel kimutatható, hogy az adott csoportban történt-e statisztikailag kimutatható fejlődés. A kontroll és a kísérleti csoportokat is összehasonlítom, független t-próbával, amely azt vizsgálja, hogy a két módszer között van-e jelentős fejlődési különbség.

A statisztikai szignifikancia mellett a hatásméretet is megmérem a Cohen-féle d értékkel. \cite{cohen1988statistical} Itt a két csoport közötti különbséget gyakorlati szempontból vizsgálom. Ez azért szükséges, mert kis minta esetén előfordul, hogy egy különbség nem tűnik statisztikailag szignifikánsnak, viszont pedagógiai vagy gyakorlati szempontból mégis figyelemreméltó tendenciát mutathat ki.

A kérdőíves válaszok elemzése leíró statisztikai módszerekkel történik. A kijelentéseket 1 és 5 között kell pontozzák, ezek alapján pedig az átlagértékből kiszámítható a motiváció, érdeklődés, használhatóság és elégedettség az alkalmazással szemben.

\section{Etikai szempontok}
A kutatás során kiemelt figyelmet szenteltem etikai szempontokra, mivel a vizsgálat során kiskorú résztvevőket mérünk fel. A részvétel önkéntes alapon történik, szülők is beleegyeztek és tájékoztatást kaptak a kutatás céljáról, adatok felhasználásáról és kezeléséről, illetve a tanár is beleegyezett, hogy az órái keretében vizsgáljak.

Az adatkezelés végig anonim módon történik. A tanulók nevei sem teszten, sem eredményekben nem szerepelnek. A résztvevők azonosítása kódokkal történik. Az így szerzett adatok kizárólag dolgozat kutatási céljára használható fel és nem alkalmas a tanulók személyes azonosítására.

A beszédkészség felmérésénél nem készülhet sem hangfelvétel, sem videófelvétel. A beszélgetés során a tanár az előre elkészített értékelési rubrika alapján közvetlenül a beszélgetés során rögzíti a pontszámokat. Ez a megoldás adatvédelmi szempontból lett alkalmazva.

A kutatás etikai szempontból fontos eleme az is, hogy az alkalmazás használata során keletkező adatok csak korlátozott célból kerüljenek felhasználásra. Az adatokat csak tanulási folyamat támogatását és kutatási elemzést szolgálják, harmadik fél számára semmi módon sem kerülhetnek továbbításra.